% DRAFT — apply during W4 restructure once DiLiGenT-MV results land.
% This subsection slots into Section 3 (method) right before \subsection{C5: matched-control deployment gate}
% to make the "what is I*?" question explicit and frame each evaluation regime
% before the C5 gate equations.

\subsection{Target signal regime}
\label{sec:target_regime}

The pipeline is parameterized by the source of the supervisory target images
\(I^*_{v,l}\) used in C1 baking, C2 candidate scoring, and C5 deployment
decisions. We evaluate the same C1--C5 mechanism under three regimes,
declared per asset:

\paragraph*{Synthetic oracle.}
For controlled assets (spot, bunny, and the Objaverse subset), the PBR
channels \(P\) are generated by a deterministic synthetic-oracle process
(\cref{sec:c1}) and \(I^*_{v,l}\) is the GGX render of those channels at
the requested view/light pair. This regime isolates UV atlas behavior
from any material-prediction error and provides a clean benchmark for
matched-control studies.

\paragraph*{Self-described oracle (procedural and SDS-init meshes).}
For procedural noisy meshes and the threestudio / Fantasia3D init shapes
in \cref{tab:transfer,tab:real_mesh}, the procedurally generated or
pipeline-attached PBR channels are treated as the oracle by construction;
\(I^*\) is rendered with the same renderer. The supervisory signal
inherits the same simplifications as the synthetic regime, but the
\emph{mesh} is no longer controlled, so this regime tests whether the
method generalizes off the synthetic oracle assets while keeping the
material model fixed.

\paragraph*{Captured imagery.}
For the DiLiGenT-MV objects (Bear, Cow, Pot2, Buddha, Reading), \(I^*\)
is the actual photographed image set, with 96 calibrated lights per view
and 20 calibrated views per object. Each captured image is loaded as a
16-bit linear PNG, normalized by the per-light RGB intensity from
\texttt{light\_intensities.txt}, masked by the silhouette (eroded
2~pixels), and excluded where any channel is saturated. No gamma or tone
mapping is applied at load time. Proposal, gate, and test partitions
are disjoint over the Cartesian product of view indices and per-view
light indices, with the partition randomized by a reported split seed.
This regime moves the supervisory signal from rendered targets to real
photometric capture, removing the synthetic-oracle assumption that
material parameters are known. The same C2 / C5 logic operates over
the captured residual without further modification.

The results in \cref{tab:diligent_mv} show that the C5 gate retains its
selective-deployment behavior under captured supervision: it accepts
when the held-out captured PSNR clears the gate threshold and rolls back
otherwise, on the same threshold set used for the synthetic regime. We
treat the captured-imagery regime as the headline validation; the
synthetic and self-described regimes serve as controlled diagnostics
(\cref{sec:matched_controls}, \cref{sec:transfer}).
